% =============================================================================
% ЛЕКЦИЯ 1
% =============================================================================

\lecture{1}{Введение}

\lecturedate{22 января 2026}

\section{Введение в программирование}

Программирование — это процесс создания компьютерных программ.

\subsection{Первая программа}

Рассмотрим классический пример программы на C++:

\begin{lstlisting}[language=C++]
#include <iostream>

int main() {
    std::cout << "Hello, World!" << std::endl;
    return 0;
}
\end{lstlisting}

хелоу ворлд так называемый
\subsection{Основные понятия}
\begin{itemize}
    \item \textbf{Переменные}: именованные области памяти для хранения данных.
    \item \textbf{Типы данных}: определяют вид данных (int, float, char и т.д.).
    \item \textbf{Операторы}: символы или комбинации символов, выполняющие операции над данными.
    \item \textbf{Условия и циклы}: конструкции для управления потоком выполнения программы.
\end{itemize}


